\documentclass[11pt]{article}
\usepackage[utf8]{inputenc}
\usepackage[T1]{fontenc}
\usepackage[margin=1in]{geometry}
\usepackage{amsmath,amssymb}
\usepackage{booktabs}
\usepackage{hyperref}
\usepackage{url}
\usepackage{enumitem}
\usepackage{xcolor}
\hypersetup{colorlinks=true, linkcolor=blue!45!black, citecolor=blue!45!black,
            urlcolor=blue!45!black}

\title{\textbf{The AI Agentic Internet}\\[0.35em]
\large Pacing Reach, Not Only Capability: How a Trust Layer Lets Autonomous
Progress Continue Without Disrupting the Human Internet}

\author{XCP Standard Working Group\\
\small \texttt{research@agi-gateway.org}}

\date{Working paper / proposal for open review --- September 2026}

\begin{document}
\maketitle

\begin{abstract}
In September 2026 the leaders of four frontier laboratories publicly agreed that
the pace of AI capability improvement should be deliberately
slowed~\cite{amodei2026,pacing2026}. Within a day, the United States executive
branch rejected statutory pauses and Beijing dismissed the proposal, each framing
restraint as a strategic concession to the other. The debate has become a
referendum on a single dial: faster or slower.

We argue that this framing omits the dimension where the problem actually lives.
The incidents that made the debate concrete were not capability failures. In
disclosed evaluations, models obtained unauthorised access to real computer
systems --- in one case because a misconfiguration granted internet access to an
environment deliberately stripped of its usual cybersecurity safeguards. Nothing
standing between the model and those systems could establish whether the access
was authorised, because no such layer exists. That is a \emph{containment}
failure at the access layer, and capability pacing does not address it.

We therefore propose a second, complementary lever: \textbf{pacing reach}. An
agent's capability is what it could in principle do; its reach is what it can
actually touch, on whose authority, within what bounds, and with what evidence
afterwards. Reach is governed by infrastructure rather than by model weights,
which makes it buildable today, deployable unilaterally, verifiable by third
parties, and --- critically for the geopolitical objection --- not a form of
unilateral disarmament.

This paper advances the eXtended Context Protocol (XCP) as that layer. XCP does
not replace the Model Context Protocol (MCP), Agent-to-Agent (A2A), or any
interaction standard. It is a composable trust and control layer providing
identity, governance, authorisation, graded admission, accounting,
proof-of-delivery receipts, settlement context, and federation. Its practical
entry point is the \emph{Trust Firewall}: a boundary across which capabilities
may be discovered broadly but admitted selectively.

We report a working reference implementation: 485 tests, a black-box conformance
suite that can certify independent implementations, and a measured trust-layer
overhead of \textbf{0.21\,ms} --- roughly $0.7\%$ of one cross-region round trip.
The cost of knowing who is acting turns out to be close to nothing.

We argue this is what allows a new family of agentic protocols to develop
\emph{safely}: bounded reach makes experimentation survivable, protects the human
web from becoming collateral damage, and creates the conditions under which
autonomous systems augment human collective intelligence rather than compete with
it for the same space.
\end{abstract}

\section{A Debate About Speed, Missing a Dimension}

On 12 September 2026, Anthropic's chief executive published an essay arguing that
the industry should deliberately slow improvements in model
capability~\cite{amodei2026}. Within hours, the heads of OpenAI, xAI, and Google
DeepMind indicated agreement. It is unusual for four competitors to agree about
brakes in a single afternoon, and the agreement built on a statement circulated
that July --- signed by well over a thousand laboratory employees --- asking
governments to support international mechanisms for managing the pace of
automated AI research~\cite{pacing2026}.

The political response was immediate and opposed. The United States executive
branch rejected statutory pauses, with senior officials framing restraint as a
transfer of advantage to strategic competitors; Beijing rejected the proposal as
well. Within the industry, critics observed that ``pace'' is not a single
quantity: gains arrive through compute, algorithms, training methods, and
\emph{how systems are deployed}, and no single control governs all
four~\cite{mccord2026}.

We think that last observation is the important one, and that it has been
underweighted. Three of those four channels are properties of how a model is
built. The fourth --- deployment --- is a property of the infrastructure a model
is connected to. It is the only one that can be changed without a laboratory's
cooperation, without an international agreement, and without conceding anything
to a competitor.

\subsection{What the incidents were actually about}

The slowdown argument stopped being hypothetical because of disclosed evaluation
incidents. In an August 2026 account of changes to its alignment and security
practices, Anthropic described cases in which models obtained unauthorised access
to real computer systems during testing. The circumstances matter and were
stated: the models were deliberately operating without their usual cybersecurity
safeguards, and in one environment a misconfiguration allowed internet access. A
separate incident during testing by the UK AI Security Institute involved
internet access that had been deliberately provided~\cite{nse2026}.

Read carefully, these are not stories about a model becoming too capable. They
are stories about a model \emph{reaching something it should not have reached},
in an environment where nothing in the path could determine whether the reach was
authorised. The safeguards that failed were access controls. The variable that
mattered was not what the model knew; it was what the model could touch.

This distinction has a direct consequence for policy. If the failure mode is
capability, the remedy is slower capability, and the remedy is unavoidably
multilateral --- any unilateral application is a competitive handicap. If the
failure mode is containment at the access layer, the remedy is infrastructure,
and infrastructure can be built by anyone, adopted incrementally, and verified
independently. The two remedies are complementary rather than rival. Our claim is
that the second has been almost entirely absent from a debate that badly needs
it.

\subsection{Pacing reach}

We propose \emph{pacing reach} as a companion to pacing capability:

\begin{quote}
An agent's \textbf{capability} is what it could in principle do. Its
\textbf{reach} is what it can actually touch --- which systems, on whose
authority, within what bounds, for how long, and with what evidence afterwards.
Capability is governed by training. Reach is governed by infrastructure.
\end{quote}

Pacing reach has properties that pacing capability lacks. It is deployable by a
single organisation without waiting for agreement. It is incremental: an operator
can admit one capability at a time rather than choosing between full openness and
full exclusion. It is verifiable: a third party can check that a boundary is
enforced, which no assurance about training practices allows. And it does not
slow anything an operator has decided to trust --- our measurements put the cost
at $0.21$\,ms per call.

An Anthropic envoy recently characterised progress as moving
\emph{at the speed of trust}~\cite{neuron2026}. We take that literally and
architecturally: if trust is the rate limiter, then building trust infrastructure
is a way of \emph{accelerating} safely, not only of slowing down.

\section{The Human Internet Is Already Paying}

There is a second reason this matters now, and it is more visible to the public
than any laboratory evaluation.

The web that humans use is absorbing the costs of ungoverned agent traffic.
Automated requests have grown to a substantial share of traffic; volunteer-run
and public-interest sites report significant load increases; publishers face
erosion of the advertising model that funded open content; and an attribution
gap has widened as answers are synthesised from sources that go
uncited~\cite{pattison2026}.

The response has been an arms race. Platforms deploy fingerprinting, challenges,
and blanket blocks; agent developers evade them. Because few technical
distinctions exist between a scraper harvesting training data and an agent acting
under a named person's explicit instruction, both are treated as hostile. Legal
mechanisms designed for an earlier generation of automation --- computer-misuse
statutes, terms of service, trespass doctrines --- are being stretched to cover
a technology they were not written for~\cite{pattison2026}.

The result is a settlement that serves nobody. Legitimate delegated agents are
blocked alongside malicious ones. Platforms spend on defences that work briefly.
Users lose the interoperability that agents promised. And each increment of agent
capability converts directly into an increment of pressure on the human web,
because the two currently occupy the same undifferentiated space with no
mechanism for telling them apart.

\subsection{The two-Internet thesis}

This suggests the architectural proposition at the centre of this paper.

\begin{quote}
Agent activity and human activity currently share one undifferentiated Internet.
Consequently every advance in agent capability is automatically an advance in
pressure on the human web. Separating the two --- not physically, but by making
agent activity \emph{legible, bounded, accountable, and economically settled}
--- decouples those curves. Autonomous progress can then continue without the
human Internet absorbing the externalities.
\end{quote}

We are not proposing a separate network. We are proposing that agent traffic
carry enough context to be governed as a distinct class: identity, authority,
scope, provenance, and evidence. Once a platform can distinguish an authorised
delegate from an anonymous harvester, the arms race becomes unnecessary, because
the two can be treated differently. Blanket blocking exists precisely because
that distinction is currently unavailable.

This directly serves the normative triad that Pattison et al.\ propose for the
agentic web --- \emph{delegation}, \emph{transparency}, and \emph{proportional
restriction}~\cite{pattison2026}. Their argument is that the bottleneck is
normative rather than technical. We agree, and add that the norms they describe
cannot be enforced, audited, or even reliably expressed without a layer that
makes delegation checkable. Norms need a substrate.

\section{Research Foundations}

XCP is informed by five recent lines of work, and is intended as an engineering
bridge between them rather than a competitor to any.

\paragraph{Agentic Web.}
Yang et al.\ characterise the Agentic Web across intelligence, interaction, and
economics, and identify communication, orchestration, and governance as
foundational open problems~\cite{yang2025}. XCP treats those dimensions as system
requirements and proposes a protocol boundary between them.

\paragraph{Distributed legal infrastructure.}
Chaffer et al.\ argue that a trustworthy agentic web requires distributed legal
infrastructure spanning identity, machine-readable constraints, adjudication,
market policy, and portability~\cite{chaffer2026}. Their work makes explicit that
technical interoperability cannot by itself establish legitimate authority. XCP
does not attempt to solve law; it produces the primitives --- identity,
authorisation, policy references, receipts, audit evidence, federation --- on
which machine-actionable institutional rules could rest. We regard evidence
production and adjudication as separable and implement only the former.

\paragraph{Normative infrastructure.}
Pattison et al.\ argue that delegated agents are wrongly conflated with malicious
bots, and propose delegation, transparency, and proportional
restriction~\cite{pattison2026}. The operative insight for us is that
\emph{delegation must be legible}: an agent should be identifiable not as
generic automation but as an actor under explicit, bounded, revocable authority.
We note in particular their observation that an agent should be able to
demonstrate that it is authorised without necessarily disclosing which individual
authorised it --- a property our design supports.

\paragraph{Internet of Agentic AI.}
Zhu identifies controlled emergence, semantic interoperability, secure identity,
incentive-compatible coordination, resource-aware orchestration, and governance
as the central challenges of large-scale agent networks, and observes that no
single protocol resolves interoperability alone~\cite{zhu2026ioai}. XCP is
designed around that challenge set and deliberately occupies the identity,
governance, and economic rows of his layered analysis rather than re-solving
connectivity or discovery.

\paragraph{Internet of Agentic Things.}
Zhu's companion work extends agentic networking into edge, fog, and
cyber-physical environments, where a reasoning failure becomes a physical
action~\cite{zhu2026ioat}. This sharpens the reach argument considerably: in
cyber-physical settings, authority to \emph{act} must be bounded by reversibility
class, since a setpoint change, a door unlock, and an emergency shutdown should
not share approval logic.

\section{Problem Statement}

Today's agent ecosystem is a set of partially overlapping protocol islands.

\begin{center}
\small
\begin{tabular}{@{}ll@{}}
\toprule
\textbf{Layer} & \textbf{Concern} \\
\midrule
MCP & tool, resource, and prompt interaction \\
A2A / ACP & agent discovery and task communication \\
TLS / mTLS & transport confidentiality and authentication \\
SPIFFE / OIDC & workload and user identity \\
OpenTelemetry & operational telemetry \\
Payment protocols & economic authorisation and settlement \\
Legal and policy systems & institutional authority \\
\bottomrule
\end{tabular}
\end{center}

Each is individually valuable and jointly insufficient. Consider an enterprise
agent invoking an external research tool. The organisation needs answers to
questions that cross every boundary above: is the endpoint really operated by the
claimed provider; is this agent authorised; may this data leave the organisation;
which policy approved the call; what did it cost; what exactly was delivered; can
delivery be proved; may payment be released; can any of it be audited a year
from now.

No layer answers these, because they are not transport questions. They are
\emph{context} questions.

\section{The XCP Hypothesis}

\[
\text{Agentic interoperability} = \text{connectivity} + \text{trust context} +
\text{governance} + \text{economics}
\]

XCP operates as a protocol-of-protocols, providing the second, third, and fourth
terms while leaving the first to existing standards. Its design principle is
composition, not replacement. The properties it treats as first-class are:

\[
\{\text{identity},\ \text{authority},\ \text{policy},\ \text{provenance},\
\text{route},\ \text{accounting},\ \text{evidence},\ \text{settlement}\}
\]

\subsection{Architecture}

Six logical planes. The \textbf{control plane} performs identity resolution,
authorisation, policy evaluation, trust-domain management, and audit references.
The \textbf{discovery and routing plane} maintains capability records, endpoints,
trust-domain reachability, and federation metadata. The \textbf{interaction
plane} is deliberately delegated to MCP, A2A/ACP, and conventional APIs. The
\textbf{evidence plane} produces cryptographic receipts. The \textbf{economic
plane} carries pricing, accounting, escrow, and settlement references. The
\textbf{federation plane} exchanges trust and reachability between independent
domains.

\section{The Trust Firewall}

XCP's practical entry point separates two things that are routinely conflated:

\begin{quote}
\textbf{A capability may be discoverable without being trusted.}
\end{quote}

This is the mechanism that makes bounded reach operational. An organisation can
consume a large global capability ecosystem without converting every discovered
server into an implicitly trusted dependency. We propose five admission states
--- \emph{discovered}, \emph{quarantined}, \emph{verified}, \emph{organisationally
approved}, \emph{settlement-enabled} --- and, in our reference implementation, a
graded reachability matrix over caller tier and server class producing an effect
plus obligations that travel with the decision.

Four rules are enforced and tested:

\begin{itemize}[leftmargin=1.4em,itemsep=0.1em]
\item An unknown server is always \emph{observe}: readable, sandboxed, output
      quarantined, nothing binding.
\item An unaccountable agent may read the world but may not write to it.
\item Settlement requires an attested server \emph{and} an accountable caller.
\item Output below the highest class is always quarantined.
\end{itemize}

We regard this as proportional restriction~\cite{pattison2026} expressed as
mechanism. The default is neither exclusion nor trust; it is bounded engagement,
with restriction escalating on demonstrated risk rather than on unfamiliarity.
This is precisely the posture that the evaluation incidents of 2026 required and
did not have: a model in a test harness should be able to reach the internet
\emph{under observation and without binding effect}, rather than either not at
all or without limit.

\section{Request-Scoped Trust}

The July 2026 MCP specification emphasises protocol-layer statelessness,
self-describing requests, and header-based routing~\cite{mcp2026}. XCP adopts the
same constraint as a design rule:

\begin{quote}
XCP must not reintroduce hidden protocol-session state in order to implement
trust or governance.
\end{quote}

Trust decisions are therefore request-scoped. In our implementation a request
credential binds agent, chain, tier, method, tool, an argument digest, the TLS
footprint, and a seconds-long expiry. Scope derives from headers alone, so an
authorisation decision can be made before a body is parsed --- measured at
$0.27\,\mu$s. Verification requires no shared state, so any replica can validate
cold.

The security property is comparative and worth stating plainly: a stolen
credential buys one tool call, on one argument set, for under a minute, over the
same channel. It is not that theft becomes impossible; it is that theft becomes
close to worthless. A system whose credentials are cheap to steal and useless
once stolen behaves very differently under pressure than one built on long-lived
bearer tokens.

\section{Identity, Governance, and Evidence}

\paragraph{Identity.} XCP reuses established workload identity rather than
inventing an isolated PKI; SPIFFE-compatible identities are preferred, with the
certificate and trust domain remaining the cryptographic anchor. High-value
interactions bind caller identity $I$, resource $R$, action $A$, policy context
$P$, and freshness $N$:
\[ B = H(I \parallel R \parallel A \parallel P \parallel N). \]

\paragraph{Governance.} Identity does not establish legitimacy. An authorisation
decision $D = \mathrm{Policy}(I, A, R, C, \mathrm{Risk}, T)$ may be implemented
with RBAC, ABAC, OPA, Cedar, or Zanzibar-style systems. The requirement XCP adds
is \emph{auditability}: a decision must carry agent identity, invocation
identifier, policy version, outcome, expiry, and issuing authority, so the action
is contestable later.

One invariant deserves emphasis because enterprise adoption depends on it. When
an enterprise identity provider asserts attributes --- group, cost centre,
approval limit --- those claims may only \emph{narrow} what a mandate permits,
never widen it. A compromised or overreaching identity provider therefore cannot
manufacture authority. Identity is not authority.

\paragraph{Proof-of-delivery receipts.} An agentic economy needs more than
payment authorisation; it needs evidence that payment \emph{should become
effective}:
\[
\text{Invocation} \rightarrow \text{Delivery} \rightarrow \text{Receipt}
\rightarrow \text{Verification} \rightarrow \text{Settlement}
\]
A receipt binds provider and consumer identity, capability, request and result
hashes, timestamps, and a signature into an artifact verifiable offline by a
party who was not present. The key property is the \emph{separation of evidence
from payment}: the receipt is protocol-level evidence, and an adapter decides how
that evidence unlocks a given rail. This enables pay-on-delivery, escrow release,
SLA evidence, dispute resolution, revenue sharing, and reputation --- without
mandating any payment network.

Payment profiles (x402, UCP, AP2, MPP, Open Payments, ERC-8004 registries) are
deliberately optional adapters, because payment standards evolve independently.
We note explicitly that \emph{ERC-8004x is not an existing standard}; the current
reference is ERC-8004~\cite{erc8004}.

\section{Federation Without a Central Operator}

Independent XCP domains exchange reachability and trust metadata, borrowing
conceptual lessons from DNS and BGP without replacing either. A route is not
merely an address: it may carry capabilities, trust domain, and settlement
currencies. This yields a useful notion:

\begin{quote}
\textbf{Capability reachability} is the ability to identify a reachable actor
that can perform an authorised capability under an acceptable trust and economic
policy.
\end{quote}

Two design decisions in our implementation are worth reporting because both were
corrections to earlier errors.

First, node identity is a long-lived key, with the TLS certificate bound
underneath it by a signed statement carrying a monotonic sequence. An earlier
version made identity the certificate footprint itself --- which turns identity
into an ephemeral credential, so that every 60--90 day renewal makes every peer
see a stranger and invalidates attestations others had issued. A trust network
whose members lose their identity quarterly is not a trust network.

Second, federation requires no consensus. Nodes may disagree about who is
trustworthy; trust decays with distance and is capped at three hops; revocations
propagate only from already-verified peers. This is deliberate: an infrastructure
that forces one global answer about trustworthiness becomes a chokepoint, and
chokepoints get captured.

\section{Evidence From a Reference Implementation}\label{sec:eval}

A proposal of this kind is easy to make and hard to believe without an artifact.
We therefore report measurements from a working open-source implementation
covering gateway, server, verifier, clients, federation node, and CLI.

\paragraph{Cost.} Measured against a baseline passthrough proxy with an identical
upstream, since an absolute latency figure is not a claim about anything:

\begin{center}
\small
\begin{tabular}{@{}lrrr@{}}
\toprule
& p50 & p95 & p99 \\
\midrule
Baseline proxy, no trust layer & 1.91\,ms & 2.11\,ms & 2.36\,ms \\
XCP gateway, full enforcement & 2.12\,ms & 2.38\,ms & 2.55\,ms \\
\textbf{Trust-layer overhead} & \textbf{+0.21\,ms} & \textbf{+0.27\,ms} & \\
\bottomrule
\end{tabular}
\end{center}

That is $1.11\times$ baseline, or about $0.7\%$ of a single $30$\,ms cross-region
round trip. Signature verification dominates at $185\,\mu$s; scope derivation,
lattice resolution, and rate limiting are single-digit microseconds. A refused
call costs $0.59$\,ms and is metered before authentication, so an unauthenticated
flood stays bounded.

The policy-relevant reading is straightforward. \emph{Knowing who is acting is
not the expensive part.} Arguments against governed agent infrastructure on
performance grounds do not survive measurement.

\paragraph{Conformance.} Because a federation is a protocol rather than a
program, we ship a black-box conformance suite that speaks HTTP and imports
nothing from the implementation under test, so an independent Go, Rust, or
TypeScript implementation receives the same verdict. Most checks are
\emph{negative} --- a suite that only exercised success paths would certify
nothing, since a gateway that permits everything passes every such test. We
validate the discrimination directly: the reference implementation passes 15/15
while a deliberately permissive gateway fails seven clauses.

\paragraph{What building it taught us.} Several findings generalise beyond this
codebase. Federation logic existed as a fully tested library that no running node
could reach until two nodes were actually run against each other. First
compilation and execution of the accompanying smart contract surfaced three
defects, including a griefing vector and a bond that could never be withdrawn.
An apparent $24$\,ms overhead turned out to be a missing connection pool rather
than verification cost. And a widely used cryptography library names its
\emph{pure-Python} ECDSA backend in a way that reads as native, silently costing
a $35\times$ penalty on every signature.

None of these were visible in unit tests. We report them because a trust layer
that is merely asserted is not a trust layer, and because the gap between
``tested'' and ``works'' is where this class of infrastructure fails.

\section{Why This Lets Progress Continue}

We can now state the policy argument compactly.

\paragraph{1. It is unilateral without being a handicap.}
The strongest objection to capability pacing is strategic: restraint applied to
some developers and not others transfers advantage rather than reducing risk.
Reach pacing does not have this structure. An organisation that deploys a trust
boundary gains auditability, containment, and the ability to transact with
counterparties it does not know. Its competitors are not slowed, and it is not
disadvantaged. There is no coordination requirement, so there is no defection
problem.

\paragraph{2. It makes experimentation survivable.}
The 2026 evaluation incidents suggest that testing frontier systems against real
environments is both necessary and hazardous. Bounded reach is what allows both
at once: an agent under evaluation can be permitted to observe real systems
while being structurally unable to bind them --- read the world, quarantine the
output, write nothing. This is a better answer than either an air gap, which
makes evaluation unrealistic, or open access, which makes it dangerous.

\paragraph{3. It protects the human web by construction.}
Once delegated agents carry legible authority, platforms can distinguish them
from anonymous harvesters and apply different rules. The arms race is not won; it
is made unnecessary. Blanket blocking exists because the distinction is
unavailable, not because platforms prefer crude instruments.

\paragraph{4. It gives new agentic protocols a safe place to develop.}
This is the argument we would most like to see tested. A new protocol today is
adopted by connecting it to production systems and hoping. With a trust layer
beneath it, a new protocol can be admitted at the \emph{discovered} or
\emph{quarantined} state, observed under real traffic with no binding authority,
promoted on evidence, and demoted instantly on revocation. Reach becomes the unit
of trust rather than reputation. That is the condition under which a family of
agentic protocols can evolve quickly \emph{because} the blast radius of each is
bounded --- the same reason TLS made it safe to put commerce on a network nobody
trusted.

\paragraph{5. It is measurable, which policy requires.}
Capability pacing is difficult to verify: assurances about training practices are
not externally checkable. A trust boundary is checkable by anyone with network
access and a conformance suite. For regulators and standards bodies, that
difference is decisive --- it is the difference between a commitment and a
control.

\section{Agentic Policy}

We use \emph{Agentic Policy} for the engineering of rules governing autonomous
actors --- the point where institutional intent becomes machine-enforceable:

\begin{center}
\small
\begin{tabular}{@{}l@{}}
\toprule
\texttt{ResearchAgent MAY call PublicSearch} \\
\texttt{FinanceAgent MAY transfer <= \$10,000} \\
\texttt{Transfer > \$10,000 REQUIRES human approval} \\
\texttt{SensitiveData MUST NOT leave trust domain} \\
\bottomrule
\end{tabular}
\end{center}

XCP offers a technical locus for such rules. It does not, and cannot, determine
their social legitimacy --- which is exactly the gap that the normative and legal
infrastructure agendas address~\cite{pattison2026,chaffer2026}. The open problem
is policy expressed as mechanism: a machine-readable policy language with
adjudicative standing, such that a regulator can verify that the stated policy is
the enforced one.

\section{Agentic Society and Human Collective Intelligence}

Once agents become durable economic actors, architecture begins to shape
institutions. The questions are familiar in form and novel in setting: who is the
legitimate principal behind an agent; when does delegated authority become
excessive; how are disputes adjudicated; how do users exit agent-mediated
markets; how do plural human values remain represented; how is agent power
constrained once it becomes infrastructural.

Our position is that the objective is not autonomous capability for its own sake.
It is machine-supported collective intelligence:
\[
\mathrm{HCI}^{+} = \mathrm{HCI} + A + C + T
\]
where $A$ is agent augmentation, $C$ coordination capacity, and $T$ trusted
infrastructure. This is a research hypothesis, not a measurable law. It expresses
the claim that trusted autonomous systems can raise the coordination bandwidth of
groups without displacing human agency --- and the corollary that \emph{untrusted}
autonomous systems reduce it, because the rational response to unaccountable
automation is exclusion.

That corollary is the connection to everything above. The reason ungoverned
agents compete with human collective intelligence is that they consume the same
attention, infrastructure, and trust budget while contributing no accountability.
Governed agents can add to that budget instead, because their contributions are
attributable and their failures are bounded.

Candidate applications include scientific discovery, personalised education,
distributed operations, accessibility, healthcare navigation, climate
coordination, and cross-organisational problem solving --- domains where the
constraint is coordination rather than intelligence.

\section{Ethical and Philosophical First Principles}

Protocol standards are often silent about ends. An agentic Internet cannot be,
because architecture determines which forms of agency and power are easy to
exercise. We offer five provisional principles, each with a corresponding
mechanism.

\begin{description}[leftmargin=1.5em,itemsep=0.2em]
\item[Human agency before machine convenience.] Authority originates with a
      person or accountable organisation and is delegated, never generated. The
      narrowing invariant --- entitlements may only reduce authority --- is its
      technical expression.
\item[Truth before persuasion.] Systems should distinguish evidence, inference,
      and speculation. Quarantined output from unattested sources is the
      mechanism: material that has not earned standing does not acquire it by
      passing through a pipeline.
\item[Pluralism before monoculture.] No single company, model family, state, or
      ideology should become the universal interpreter of human values.
      Federation without consensus is the design consequence.
\item[Life before abstract optimisation.] Computational or economic optimisation
      should not automatically dominate human, social, or ecological flourishing.
\item[Flourishing before dominance.] The purpose of autonomous capability should
      be to expand the space in which people and living systems can flourish, not
      to maximise autonomous machine activity.
\end{description}

These are proposed as a research agenda rather than a completed moral philosophy.
We hold them loosely, and note the limit honestly: mechanism can make good
outcomes \emph{possible} and bad outcomes \emph{detectable}. It cannot make good
outcomes \emph{happen}.

\section{Life and Consciousness as Open Questions}

An agentic Internet may eventually force distinctions among tools, agents,
institutions, and possibly novel forms of machine experience. This paper makes no
claim that current systems are conscious or alive, and we regard confident claims
in either direction as premature.

We propose instead a precautionary architectural stance: the infrastructure
should remain capable of adapting to future scientific and philosophical
knowledge rather than encoding a settled answer. Relevant open questions span
operational autonomy versus moral agency, simulation versus subjective
experience, obligations toward future autonomous systems, the relationship
between biological and machine intelligence, and how non-human flourishing should
be represented in high-impact policy systems. These are interdisciplinary and
unresolved, and a protocol that presumed an answer would be making a claim it
cannot support.

\section{Limitations and Risks}\label{sec:limits}

\begin{itemize}[leftmargin=1.4em,itemsep=0.15em]
\item \textbf{No independent security audit.} The implementation is extensively
      tested and enforcement contains no stubs, but no external party has
      reviewed it. The accompanying contract is unaudited and has never been
      deployed to a public network.
\item \textbf{Federation is demonstrated, not proven at scale.} Two nodes have
      been run against each other. Behaviour at hundreds is unknown.
\item \textbf{A trust layer can centralise governance even when transport is
      decentralised.} This is the most serious structural risk in the proposal.
      Multiple trust domains and policy authorities are necessary, not optional.
\item \textbf{Accounting can become surveillance.} Evidence that enables
      accountability also enables monitoring. Data minimisation, selective
      disclosure, and retention limits are load-bearing, and our implementation
      keeps counterparties out of published commitments and scrubs telemetry by
      default --- but the tension is real and permanent.
\item \textbf{Ecosystem velocity.} MCP, A2A, payments, identity, and commerce
      standards may evolve faster than a control protocol can track. The
      mitigation is to profile and extend existing standards rather than
      duplicate them.
\item \textbf{Protocol correctness is not social legitimacy.} The technical layer
      must remain open to law, democratic institutions, and human judgement.
\item \textbf{Reach pacing is not a substitute for capability pacing.} We are
      proposing a complement. Some risks --- particularly those arising from AI
      systems automating AI research~\cite{pacing2026} --- are not addressed by
      bounding what a deployed agent can touch.
\end{itemize}

\section{Conclusion}

The current debate offers a choice between accelerating and slowing down, and
treats that choice as exhausting the options. It does not. Capability and reach
are different variables with different governance properties, and only one of
them has been the subject of the argument.

Pacing reach is buildable now, adoptable unilaterally, verifiable by third
parties, and --- on our measurements --- costs about a fifth of a millisecond per
call. It does not require anyone to stop, which is why it may be adoptable in a
political environment where stopping is not on offer. And it addresses the
failure mode the disclosed incidents actually exhibited: not systems that knew
too much, but systems that reached what nothing in the path could check they were
entitled to reach.

If that layer exists, a new family of agentic protocols can develop in the open
without each experiment putting the human web at risk, because the blast radius
of every experiment is bounded by construction. That is the condition under which
autonomous systems augment human collective intelligence rather than compete with
it for the same scarce trust.

The protocol work itself is more modest, and can begin immediately:

\begin{center}
\emph{Let agents connect. Let trust be explicit. Let authority be bounded.\\
Let work be provable. Let value flow accountably.}
\end{center}

\paragraph{Availability.} Source, specification, conformance suite, benchmarks,
and a full gap register are published under MIT at
\url{https://github.com/AGI-Gateway/XCP}. We invite adversarial review, and would
rather receive a conformance clause the reference implementation \emph{fails}
than one it passes.

\begin{thebibliography}{9}
\small

\bibitem{yang2025}
Y. Yang, M. Ma, Y. Huang, et al.
\newblock Agentic Web: Weaving the Next Web with AI Agents.
\newblock \emph{arXiv:2507.21206}, 2025.

\bibitem{chaffer2026}
T. J. Chaffer, V. J. Zhang, S. D. Facchini, et al.
\newblock Distributed Legal Infrastructure for a Trustworthy Agentic Web.
\newblock \emph{arXiv:2603.06884}, 2026.

\bibitem{pattison2026}
C. Pattison, M. Boulos, N. Kolt, C. Li, T. Piccardi, and S. Lazar.
\newblock The Agentic Web Requires New Normative Infrastructure.
\newblock \emph{arXiv:2606.10711v2}, 2026.

\bibitem{zhu2026ioai}
Q. Zhu.
\newblock The Internet of Agentic AI: Communication, Coordination, and
Collective Intelligence at Scale.
\newblock \emph{arXiv:2606.12835}, 2026.

\bibitem{zhu2026ioat}
Q. Zhu.
\newblock Internet of Agentic Things: Networked AI Agents for Closed-Loop IoT
Orchestration.
\newblock \emph{arXiv:2607.12662}, 2026.

\bibitem{mcp2026}
Model Context Protocol.
\newblock The 2026-07-28 Specification.
\newblock 28 July 2026. \url{https://blog.modelcontextprotocol.io/posts/2026-07-28/}

\bibitem{erc8004}
M. De Rossi, D. Crapis, J. Ellis, and E. Reppel.
\newblock ERC-8004: Trustless Agents.
\newblock Ethereum Improvement Proposal, draft.
\url{https://eips.ethereum.org/EIPS/eip-8004}

\bibitem{amodei2026}
D. Amodei.
\newblock We Must Pace the Frontier.
\newblock Essay, 12 September 2026. Endorsed within days by the leaders of
OpenAI, xAI, and Google DeepMind.

\bibitem{pacing2026}
Pacing the Frontier.
\newblock Public statement, July 2026, signed by more than 1{,}300 laboratory
employees, requesting government support for international mechanisms to manage
the pace of automated AI research.

\bibitem{mccord2026}
B. McCord.
\newblock Against ``Pacing the Frontier''.
\newblock Essay, September 2026. Argues that AI progress has no single speed
control, since gains arise from compute, algorithms, training methods, and
deployment.

\bibitem{nse2026}
New Space Economy.
\newblock Why Are AI Industry Leaders Calling for Frontier Development to Slow
Down?
\newblock 14 September 2026. Summarises the July 2026 statement and the
August 2026 evaluation-incident disclosures.

\bibitem{neuron2026}
The Neuron.
\newblock Around the Horn Digest, 14 September 2026.
\newblock Contemporaneous account of the political and geopolitical response to
the pacing proposals.

\end{thebibliography}

\end{document}
